\subsection[$\real_z$-actions: noncommutative torus bundles]{$\boldsymbol{\real_z}$-actions: noncommutative torus bundles}\label{sec:Rzgen}

Let us now come to the torsion summands $\Z_{m_i}$ in the
$\Z$-module decomposition \eqref{eq:H1gendecomp}. Pick a~non-trivial
elementary divisor $m_i>1$ for some $i\in\{1,\dots,r\}$. The
corresponding homology generator $\tilde e_i$ is constructed as a~$\Z$-linear combination~\eqref{eq:tildeei} of the generators~\eqref{eq:hatei}. It defines a~fixed
element~$z_0$ in the image of $\varphi_{x_0}-\rm{id}_{\R^{d-1}}$ and a~corresponding one-dimensional subgroup of~$\sfG$ given by
\[
\R_{z_0} := \R (z_0,0) .
\]
We further assume that
\[
\Z_{ z_0} := \R_{ z_0}\cap\Lambda_\sfG
\]
is a lattice in $\R_{z_0}$.
The choice of $z_0\in\R^{d-1}$ is not unique, and
any change of basis of $\Z^{d-1}$, represented by a matrix
$B\in\sfGL(d-1,\Z)$, defines an equally good element $B\cdot
z_0\in\R^{d-1}$ as long as $B\cdot
z_0\in{\rm im}\big(\varphi_{x_0}-\Id_{\R^{d-1}}\big)$. We write $\proj_{z_0}$ for the
linear projection of $\R^{d-1}$ to $\R_{z_0}$, and denote by $\langle
z_0,z\rangle\in\R$ the component of $z\in\R^{d-1}$ in $\R_{z_0}$,
i.e.,~$\proj_{z_0}\cdot z = \langle z_0,z\rangle \, z_0$.

In the basis $\hat e_i$, the lattice of $\sfG$ is given by
\[
\Lambda_\sfG=\Z^{d-1}\rtimes_{\hat\varphi|_{x_0\,\Z}} x_0 \Z ,
\]
where $\hat\varphi_x:=\Sigma^{-1}\,\varphi_x\,\Sigma$ for all
$x\in\R$. Then the fibres of the underlying Mostow bundle are `square' tori
$\torus^{d-1}$ with unit periodicities $ z\sim z+\vec
e_i$ for $i=1,\dots,d-1$, where as
before $\vec e_i$ denotes the standard basis of
$\R^{d-1}$. The action of the
subgroup $\R_{ z_0}$ on
elements $(z,x)\in\sfG$ by right multiplication is given by
\begin{gather}\label{eq:Rz0action}
(z,x) (\zeta\,z_0,0) = \big(z+\zeta \Sigma^{-1} \varphi_x\cdot
z_0,x\big)
\end{gather}
for $\zeta\in\R$, where $\Sigma^{-1} \varphi_x\cdot z_0$ lies in the
image of the relation matrix ${\tt A}=\ttM-\unit_{d-1}$.

Our principal tool to compute the $C^*$-algebraic
T-dual for such an action of $\R$ in the image of $\varphi_{x_0}-\Id_{\R^{d-1}}$, which we collectively refer to as
$\R_z$-actions, will be Green's theorem in the form~\eqref{scheme}:
\begin{gather}\label{Mostowscheme}
C(\torus_{\Lambda_\sfG})\rtimes_\rt \real_{ z_0} \ \mbf\sim_{\text{\tiny M}} \
C_0(\sfG/\real_{ z_0})\rtimes_\lt \Lambda_\sfG .
\end{gather}
By appealing to
Theorem~\ref{thm:crossedfibres}, we may apply \eqref{Mostowscheme}
fibrewise. For fixed $x\in\R$, the fibre $\sfG_x$ of the semi-direct
product $\sfG=\R^{d-1}\rtimes_{\hat\varphi}\R$ is the subgroup~$\R^{d-1}$, and the corresponding fibre of
the solvmanifold $\torus_{\Lambda_\sfG}$ over $x\in\R/x_0\,\Z$ is the
torus $\torus^{d-1}=\R^{d-1}/\Z^{d-1}$. The Morita
equivalence~\eqref{Mostowscheme} is $C(\torus)$-linear and the fibres of the corresponding
T-dual $C(\torus)$-algebra are given by the fibrewise Morita
equivalence
\begin{gather}\label{eq:Mostowschemefiber}
\big(C(\torus_{\Lambda_\sfG})\rtimes_\rt\real_{ z_0}\big)_x \simeq
C\big(\R^{d-1}/\Z^{d-1}\big)\rtimes_{\rt^x}\real_{ z_0} \
\mbf\sim_{\text{\tiny M}} \
C_0\big(\real^{d-1}/\R_{ z_0}\big)\rtimes_{\lt^x}\Z^{d-1} .
\end{gather}
The action of the subgroup $\Z^{d-1}\subset\Lambda_\sfG$ on the
coset space $\R^{d-1}/\R_{ z_0}$ is induced by left multiplication in the
group $\sfG$.
After a basis transformation, we can decompose the discrete group $\Z^{d-2}$ into a direct sum
\[
\Z^{d-1} \simeq \Z_v^{d-2}\oplus\Z_{ z_0} ,
\]
where $\Z_v^{d-1} = (\unit_{d-1}-\proj_{z_0})\cdot\Z^{d-1}$.

Let
\[
\mathbb{F}_* := \big\{x\in\R \, \big| \, \big\langle
z_0,\Sigma^{-1} \varphi_x\cdot z_0\big\rangle=0\big\} .
\]
\begin{Lemma}\label{lem:ellfibercom}
Over any $x\in \mathbb{F}_*$, the fiber
$\big(C(\torus_{\Lambda_\sfG})\rtimes_\rt\real_{ z_0}\big)_x$
is Morita equivalent to the
commutative $C^*$-algebra \smash{$C\big(\torus^{d-1}\big)$}.
\end{Lemma}
\begin{proof}
If $x\in\mathbb{F}_*$, then $w_0:=\Sigma^{-1}\,\varphi_x\cdot z_0$ only
shifts the corresponding component of $ z$ in
$(\unit_{d-1}-\proj_{ z_0})\cdot\R^{d-1}$ in the $\R_{z}$-action
\eqref{eq:Rz0action}. In this case any
element $(z,x)\in\sfG$ may be factorized as
$$
(z,x) = \big(\proj_{ z_0}\cdot
z+(\unit_{d-1}-\proj_{z_0}-\proj_{w_0})\cdot z\,,\,x\big) \,
\big(\langle w_0,z\rangle\, z_0 \,,\,0\big) \ .
$$
Hence the coset space $\R^{d-1}/\R_{ z_0}\simeq\R^{d-2}$ can be
parameterized by $\proj_{ z_0}\cdot
z+\big(\unit_{d-1}-\proj_{z_0}-\proj_{w_0}\big)\cdot z$ with $z\in\R^{d-1}$,
which we decompose correspondingly into a direct product
$\R_{z_0}\times\R^{d-3}$ (with the second factor absent for $d=3$). By \eqref{eq:Mostowtorus} the discrete group
$\Z_v^{d-2}$ acts trivially on the line $\R_{z_0}$ and by translations
on $\R^{d-3}$,
while~$\Z_{ z}$ acts by translations on $\R_{z_0}$ and trivially
on $\R^{d-3}$. Then the crossed product on the right-hand side of~\eqref{eq:Mostowschemefiber} may be unravelled to get
\begin{align*}
C_0\big(\real^{d-1}/\R_{ z_0}\big)\rtimes_{\lt^x}\Z^{d-1}
&\simeq C_0\big(\R_{
 z_0}\times\R^{d-3}\big)\rtimes_{\lt^x}\big(\Z_v^{d-2}\times\Z_{ z_0}\big) \\
&\simeq
\big[\big(C_0(\R_{z_0})\otimes C_0\big(\R^{d-3}\big)\big)\rtimes_{\Id\otimes\lt^x}\Z^{d-2}_v\big]\rtimes_{\beta^x}\Z_{z_0} \\
& \ \mbf\sim_{\text{\tiny M}} \ \big(C_0(\R_{ z_0})\otimes C^*(\Z)\otimes C\big(\torus^{d-3}\big)\big)\rtimes_{\lt^x\otimes\Id\otimes\Id}\Z_{ z_0}
 \\
&\simeq \big(C_0(\R_{ z_0})\rtimes_{\lt^x}\Z_{ z_0}\big)\otimes
C^*(\Z)\otimes C\big(\torus^{d-3}\big)
\nonumber\\
& \ \mbf\sim_{\text{\tiny M}} \ C(\torus)\otimes C^*(\Z)\otimes C\big(\torus^{d-3}\big) \simeq C\big(\torus^{d-1}\big) .
\end{align*}
In the second line we applied
Theorem~\ref{cpsp3}. In the third line we used
Example~\ref{ex:Moritatori} together with the fact that the
homomorphism $\sigma_{\Z_{ z_0}}$ from \eqref{sigmaI} is trivial since the
groups $\Z_{ z_0}$ and $\Z_v^{d-2}$ are discrete, and so the action of $\Z_{ z_0}$ on the crossed
product is induced by left multiplication on $\R_{ z_0}$, the trivial
action on $\R^{d-3}$, and the
trivial action on $\Z_v^{d-2}$. In the fifth line we used
Example~\ref{ex:Moritatori} again.
\end{proof}

The central result of this paper is
\begin{Theorem}
\label{thm:NCtorusbundlegen}
The $C^*$-algebraic T-dual of any almost abelian
solvmanifold $\torus_{\Lambda_\sfG}$ with respect to an $\R_z$-action
is Morita equivalent to a $C^*$-algebra bundle
of noncommutative tori over the circle $\torus$:
\begin{gather*}%\label{eq:algbundlegen}
C(\torus_{\Lambda_\sfG})\rtimes_\rt\real_{ z_0}
 \ \mbf\sim_{\text{\tiny M}} \ \mcoprod
\torus^{d-1}_{\vec\theta_{z_0}(x)} ,
\end{gather*}
where the noncommutativity parameters $\vec\theta_{z_0}(x)\in\R^{d-2}$ are given
by
\begin{gather}\label{eq:thetaz0}
\vec\theta_{z_0}(x) = \begin{cases}
 0 & \mbox{for} \ x \in \F_* , \\
 \dfrac{(\unit_{d-1}-\proj_{ z_0}) \Sigma^{-1} \varphi_{x_0 x}\cdot
 z_0}{\big\langle{ z_0},\Sigma^{-1} \varphi_{x_0 x}\cdot z_0\big\rangle}
 & \mbox{for} \ x\in\R\setminus\F_* . \end{cases}
\end{gather}
\end{Theorem}
\begin{proof}
That the fibres over $x\in\F_*$ are just ordinary tori $\torus^{d-1}$
is established by Lemma~\ref{lem:ellfibercom}, so we may assume that $x\in\R\setminus\F_*$. Then a simple
calculation shows that any element $(z,x)\in\sfG$ can be factorized as
\[
(z,x) = (v,x) \left(\frac{\proj_{ z_0}\cdot z}{\big\langle z_0,\Sigma^{-1} \varphi_x\cdot z_0\big\rangle} , 0\right) ,
\]
where
\begin{gather}\label{eq:Mostowcoset}
v:=\big(\unit_{d-1}-\proj_{ z_0}\big)\cdot z - \frac{\langle{ z_0},
 z\rangle}{\big\langle{ z_0},\Sigma^{-1} \varphi_x\cdot z_0\big\rangle}
\big(\unit_{d-1}-\proj_{ z_0}\big) \Sigma^{-1} \varphi_x\cdot z_0 .
\end{gather}
Thus the coset space $\R_{x,v}^{d-2}:=\R^{d-1}/\R_{ z_0}$ may be
parameterized by the coordinates $v\in\R^{d-2}$ over any
$x\in\R\setminus\F_*$, and we explicitly retain the fibre index in the
notation for convenience.

We now need to unravel the crossed product
on the right-hand side of~\eqref{eq:Mostowschemefiber}.
From~\eqref{eq:Mostowtorus} and~\eqref{eq:Mostowcoset} it follows that the action of elements
$(\gamma_v,\gamma_{ z_0})\in\Z^{d-1}\simeq\Z_v^{d-2}\oplus\Z_{
 z_0}$ on the coset is given by
\begin{gather}\label{eq:Zcosetactionv}
(\gamma_v,0,0)\cdot (v,x) = (v+\gamma_v,x) , \\
(0,\gamma_{ z_0},0)\cdot (v,x) = \left(v-\frac{
 \gamma_{ z_0}}{\big\langle{ z_0},\Sigma^{-1} \varphi_x\cdot z_0\big\rangle}
\big(\unit_{d-1}-\proj_{ z_0}\big) \Sigma^{-1} \varphi_x\cdot
 z_0 , x\right) .
\label{eq:Zcosetactionz}\end{gather}
Applying
Theorem~\ref{cpsp3} as in the proof of Lemma~\ref{lem:ellfibercom} gives
\begin{gather*}%\label{eq:MostowRHS}
C\big(\R^{d-1}/\Z^{d-1}\big)\rtimes_{\rt^x}\real_{ z_0}
\ \mbf\sim_{\text{\tiny M}} \
\big(C_0\big(\real^{d-2}_{x,v}\big)\rtimes_{\lt^x}\Z^{d-2}_{v}\big) \rtimes_{\beta^x} \Z_{ z_0} ,
\end{gather*}
where the action of $\Z^{d-2}_v$ on the coset $\R^{d-2}_{x,v}$ is
given by~\eqref{eq:Zcosetactionv}, while the action of $\Z_{ z_0}$ on the crossed product
$C_0\big(\real^{d-2}_{x,v}\big)\rtimes_{\lt^x}\Z^{d-2}_{v}$ is induced by the
action on $\real^{d-2}_{x,v}$ given in~\eqref{eq:Zcosetactionz} and the trivial action on~$\Z^{d-2}_{v}$.

Next, Example~\ref{ex:Moritatori} yields
\begin{gather}\label{eq:NMoritav}
C_0\big(\real_{x,v}^{d-2}\big)\rtimes_{\lt^x} \Z_{v}^{d-2} \ \mbf\sim_{\text{\tiny M}} \
C\big(\torus^{d-2}\big) .
\end{gather}
Since the homomorphism
$\sigma_{\Z_{ z_0}}\colon \Z_{ z_0}\to\real^+$ is trivial, it is not difficult to see
that the Morita equivalence bimodule $\Mcal$ implementing this
equivalence is $\Z_{ z_0}$-equivariant (in the sense of Theorem~\ref{equivMorita}):
the $\Z_{ z_0}$-action
$U\colon \Z_{ z_0}\rightarrow \Aut(\Mcal)$ is given by $U_{\gamma_v}(\xi)(v)=\xi(v+\gamma_v)$, for all $\gamma_v\in
\Z_{ z_0}$, $\xi\in \Mcal$ and $v\in \real_{x,v}^{d-2} $.
Applying Theorem \ref{equivMorita} we conclude that the
Morita equivalence \eqref{eq:NMoritav} induces the Morita equivalence
\begin{gather}\label{eq:fibreNCMostow}
\big(C_0\big(\real^{d-2}_{x,v}\big)\rtimes_{\lt^x}\Z^{d-2}_{v}\big) \rtimes_{\beta^x}
\Z_{ z_0} \ \mbf\sim_{\text{\tiny M}} \
C\big(\torus^{d-2}\big)\rtimes_{\lt^x}\Z_{ z_0} ,
\end{gather}
where the action
of the group $\Z_{ z_0}$ on $C\big(\torus^{d-2}\big)$ is the pullback of the action on
$\torus^{d-2}=\real_{x,v}^{d-2}/\Z_{v}^{d-2}$ induced
from~\eqref{eq:Zcosetactionz}.

After rescaling $x$ by $x_0$
to give it unit period, this shows that the
corresponding algebra of functions on the fiber at $\e^{2\pi\ii x}$ is that
of a noncommutative $d{-}1$-torus $\torus_{\vec\theta_{
 z_0}(x)}^{d-1}$ with noncommutativity parameter
$\vec\theta_{z_0}(x)$ given by~\eqref{eq:thetaz0}, see Example~\ref{ex:NCtorid}.
Thus the $C^*$-algebraic T-dual
$C(\torus_{\Lambda_\sfG})\rtimes_\rt\real_{ z_0}$ is a $C^*$-algebra bundle
of noncommutative tori $\A_{\vec\theta_{
 z_0}(x)}=\torus^{d-1}_{\vec\theta_{ z_0}(x)}$ over the
circle $\torus=\big\{\e^{2\pi\ii x}\, | \, x\in \real/\Z\big\}$.
\end{proof}

What becomes of the non-trivial monodromy \eqref{eq:Mostowtorus} of
the $\torus^{d-1}$ fibers parameterized by~$z$ in the original Mostow
bundle? To answer this question, we write the monodromy matrix
$\ttM=(m_{ij})\in\sfSL(d-1,\Z)$ in the block form
\begin{gather}\label{eq:ttMblock}
\ttM = \left(\begin{matrix} \ttM_{|d-2} & \vec m \\
\vec m' & m_{d-1 \, d-1}
\end{matrix}\right) ,
\end{gather}
where $\ttM_{|d-2}=(m_{ij})_{1\leq i,j\leq d-2}$, while $\vec m=(m_{i \,
 d-1})_{i=1,\dots,d-2}$ and $\vec m' = (m_{d-1 \,
 i})_{i=1,\dots,d-2}$ are respectively column and row vectors in $\Z^{d-2}$. Below
we denote the usual Euclidean inner product on $\R^{d-2}$ by
$\langle\,\cdot\,,\,\cdot\,\rangle$.
\begin{Proposition}\label{prop:monodromyprop}
The noncommutativity parameter
$\vec\theta_{\vec e_{d-1}}(x)$ from~\eqref{eq:thetaz0} varies with a
change of coset representative $x\in\R/\Z$ according to its
$\sfSL(d-1,\Z)$ orbit under the action of the monodromy matrix
\eqref{eq:ttMblock} by $(d{-}2)$-dimensional linear fractional transformations
\begin{gather}\label{eq:thetaMoritagen}
\vec\theta_{\vec e_{d-1}}(x+1) = \ttM\big[\vec\theta_{\vec
 e_{d-1}}(x)\big] :=
\frac{\ttM_{|d-2}\cdot\vec\theta_{\vec e_{d-1}}(x) + \vec m}
{\big\langle \vec m',\vec\theta_{\vec e_{d-1}}(x)\big\rangle + m_{d-1 \,
 d-1}} .
\end{gather}
Under a change of basis of $\Z^{d-1}$ given
by a matrix $B\in\sfGL(d-1,\Z)$, the corresponding noncommutativity
parameter varies according to
\[
\vec\theta_{B\cdot\vec e_{d-1}}(x+1) = \big(B^{\rm
 t} \ttM B\big)\big[\vec\theta_{B\cdot\vec e_{d-1}}(x)\big] .
\]
\end{Proposition}
\begin{proof}
The key feature stems from the definition \eqref{eq:monodromydef} of the monodromy matrix
$\ttM$:
\[
\Sigma^{-1}\,\varphi_{x+x_0} = \Sigma^{-1} \varphi_{x_0} \varphi_x =
\ttM \Sigma^{-1} \varphi_x .
\]
The statements then follow from straightforward calculations in components.
\end{proof}

\begin{Remark}
\label{rem:MoritaKK}
The right-hand side of \eqref{eq:thetaMoritagen} is an example of a
linear fractional transformation in higher dimensions, known from
complex analysis, see, e.g.,~\cite{Cowen2000}. We can show that it defines a Morita equivalence of noncommutative $d{-}1$-tori from
Example~\ref{ex:NCTdMorita}. For this, we introduce the skew-symmetric
matrix $\Theta$ corresponding to the vector $\vec\theta\in\R^{d-2}$ as
in Example~\ref{ex:NCtorid}:
\[
\Theta = \left(
\begin{matrix}
\mbf 0_{d-2} & \vec\theta \\ -\vec\theta\,^{\rm t} & 0
\end{matrix}
\right) .
\]
We denote by $\mbf g_{\ttM}$ the element of ${\sf SO}(d-1,d-1;\Z)$
corresponding to the monodromy matrix $\ttM\in\sfSL(d-1,\Z)$:
\[
\mbf g_{\ttM} = \left(
\begin{matrix}
\ttM & \mbf 0_{d-1} \\ \mbf 0_{d-1} & \big(\ttM^{\rm t}\big)^{-1}
\end{matrix}
\right)
\in {\sf SO}(d-1,d-1;\Z) .
\]
Introduce matrices $\mbf T_{d-1}$ of order $2$ with determinant $-1$
by
\[
\mbf T_{d-1} = \left(
\begin{matrix}
\unit_{d-1}-{\tt E}_{d-1} & {\tt E}_{d-1} \\ {\tt E}_{d-1} &
\unit_{d-1}-{\tt E}_{d-1}
\end{matrix}
\right) \in {\sf O}(d-1,d-1;\Z) ,
\]
where ${\tt E}_{d-1}$ is the matrix unit whose only
non-zero element is $({\tt E}_{d-1})_{d-1\,d-1} = 1$.
Finally, define the element $\mbf M\in{\sf SO}(d-1,d-1;\Z)$ by
\[
\mbf M = \mbf T_{d-1} \mbf g_{\ttM} \mbf T_{d-1} .
\]
Using the adjugate formula
for matrix inverses, a straightforward if tedious calculation then shows that the
corresponding ${\sf SO}(d-1,d-1;\Z)$ orbit of $\Theta$ reproduces the
$(d{-}2)$-dimensional linear fractional transformation of
\eqref{eq:thetaMoritagen}:
\[
\mbf M[\Theta] = \left(
\begin{matrix}
\mbf 0_{d-2} & \ttM\big[\vec\theta\, \big] \\
-\ttM\big[\vec\theta\, \big]^{\rm t} & 0
\end{matrix}
\right) \qquad \mbox{with} \quad \ttM\big[\vec\theta\, \big] = \frac{\ttM_{|d-2}\cdot\vec\theta + \vec m}
{\big\langle \vec m',\vec\theta\, \big\rangle + m_{d-1 \,
 d-1}} .
\]
It follows that, while the fibre noncommutative
tori of Theorem~\ref{thm:NCtorusbundlegen} are not generally identical under a change of
representative $x\in\R/\Z$, by Example~\ref{ex:NCTdMorita}
they are always Morita
equivalent:
\[
\torus^{d-1}_{\vec\theta_{\vec e_{d-1}}(x+1)} \ \mbf\sim_{\text{\tiny
 M}} \ \torus^{d-1}_{\vec\theta_{\vec e_{d-1}(x)}} .
\]
Recalling that our formulation of topological T-duality takes place in the additive category~${\CCK\hspace{-1mm}\CCK}$ from
Section~\ref{sec:TdualityKK}, this has a natural interpretation: The
non-trivial monodromy in the automorphism group of the~\smash{$\torus^{d-1}$} fibres of the
twisted torus $\torus_{\Lambda_\sfG}$ is
manifested as a (generally non-trivial) isomorphism in the Morita automorphism
group of the fibre noncommutative tori in~${\CCK\hspace{-1mm}\CCK}$.
\end{Remark}

